% GENERATED FILE. Edit canonical lessons, then run npm run generate:books.
% canonical-source-hash: fnv1a64:c8d1ddeca3192f29
% canonical-lessons: JA-C432-nayamu, JA-C432-kuyokuyo-suru, JA-C432-koukai-suru, JA-C432-munashii, JA-C432-akirameru

\chapter{Worry and Regret}
\label{ch:ja-worry-and-regret}
\hlchaptermodality{432}

\begin{chapteropening}
\textbf{By the end of this chapter:} \emph{I can say what worries me, what I regret and what I gave up.}

Complete the last lesson of chapter 432: I can say what worries me, what I regret and what I gave up.
\end{chapteropening}

\section[\texorpdfstring{\ja{なやむ} (nayamu)}{なやむ (nayamu)}]{\ja{なやむ} (nayamu) — to be troubled; to agonize}

\label{lesson:JA-C432-nayamu}



\begin{quote}
Before the new one: say the Japanese for to surprise someone, then the Japanese for to be on edge.
\end{quote}

\subsection*{You'll want to know: \ja{なやむ}}
\textbf{\ja{なやむ}} — \emph{nayamu} — \textquotedblleft{}to be troubled; to agonize\textquotedblright{}.

\begin{quote}
\textbf{\ja{あには} \ja{しごとで} \ja{なやんでいます。}} — \emph{ani wa shigoto de nayande imasu} — \textquotedblleft{}My elder brother is troubled about his work.\textquotedblright{}
\end{quote}

Polite \textbf{\ja{なやみます}}, \emph{nayamimasu}; a godan verb. Past (polite): \textbf{\ja{ともだちが} \ja{なやみました。}}, \emph{tomodachi ga nayamimashita}.

\begin{cousinweb}
\textbf{\ja{なやむ}} (\emph{nayamu}) is a native Japanese word (wago), inherited rather than borrowed from Chinese; its kanji spelling was attached to it later (JMdict).
\end{cousinweb}

\begin{grammarlens}[title={What you've built}]
One more word for describing what you lived and felt.
\end{grammarlens}

\subsection*{Your turn}
\begin{itemize}
\raggedright
  \item \emph{Say it:} \emph{nayamu}
  \item \emph{Say it:} \emph{nayamu}, once more
  \item \emph{Say it:} \emph{ani wa shigoto de nayande imasu} — My elder brother is troubled about his work.
  \item \emph{Recall:} say the Japanese for to surprise someone, then the Japanese for to be on edge, then say \emph{nayamu} again
\end{itemize}

\begin{tcolorbox}[breakable,colback=teal!4,colframe=teal!35!black,title={Before you move on}]
What does \ja{なやむ} mean? (\textquotedblleft{}To be troubled; to agonize\textquotedblright{}.) Say it once more.
\end{tcolorbox}
\section[\texorpdfstring{\ja{くよくよする} (kuyokuyo suru)}{くよくよする (kuyokuyo suru)}]{\ja{くよくよする} (kuyokuyo suru) — to brood; to fret}

\label{lesson:JA-C432-kuyokuyo-suru}



\begin{quote}
Before the new one: say the Japanese for to be on edge, then the Japanese for to be troubled.
\end{quote}

\subsection*{You'll want to know: \ja{くよくよする}}
\textbf{\ja{くよくよする}} — \emph{kuyokuyo suru} — \textquotedblleft{}to brood; to fret\textquotedblright{}.

\begin{quote}
\textbf{\ja{ちいさい} \ja{ことで} \ja{くよくよしないで} \ja{ください。}} — \emph{chiisai koto de kuyokuyo shinaide kudasai} — \textquotedblleft{}Please do not brood over small things.\textquotedblright{}
\end{quote}

Polite \textbf{\ja{くよくよします}}, \emph{kuyokuyoshimasu}; a \ja{する} verb. Past (polite): \textbf{\ja{かれが} \ja{くよくよしました。}}, \emph{kare ga kuyokuyo shimashita}.

\begin{cousinweb}
\textbf{\ja{くよくよする}} (\emph{kuyokuyo suru}) is a native Japanese mimetic word, \emph{kuyokuyo}, inherited rather than borrowed from Chinese, made a verb with \emph{suru}, \textquotedblleft{}to do\textquotedblright{}; its older origin is not settled (JMdict).
\end{cousinweb}

\begin{grammarlens}[title={What you've built}]
One more word for describing what you lived and felt.
\end{grammarlens}

\subsection*{Your turn}
\begin{itemize}
\raggedright
  \item \emph{Say it:} \emph{kuyokuyo suru}
  \item \emph{Say it:} \emph{kuyokuyo suru}, once more
  \item \emph{Say it:} \emph{chiisai koto de kuyokuyo shinaide kudasai} — Please do not brood over small things.
  \item \emph{Recall:} say the Japanese for to be on edge, then the Japanese for to be troubled, then say \emph{kuyokuyo suru} again
\end{itemize}

\begin{tcolorbox}[breakable,colback=teal!4,colframe=teal!35!black,title={Before you move on}]
What does \ja{くよくよする} mean? (\textquotedblleft{}To brood; to fret\textquotedblright{}.) Say it once more.
\end{tcolorbox}
\section[\texorpdfstring{\ja{こうかいする} (kōkai suru)}{こうかいする (kōkai suru)}]{\ja{こうかいする} (kōkai suru) — to regret}

\label{lesson:JA-C432-koukai-suru}



\begin{quote}
Before the new one: say the Japanese for to be troubled, then the Japanese for to brood.
\end{quote}

\subsection*{You'll want to know: \ja{こうかいする}}
\textbf{\ja{こうかいする}} — \emph{kōkai suru} — \textquotedblleft{}to regret\textquotedblright{}.

\begin{quote}
\textbf{\ja{あの} \ja{ときの} \ja{ことを} \ja{こうかいしています。}} — \emph{ano toki no koto o kōkai shite imasu} — \textquotedblleft{}I regret what happened that time.\textquotedblright{}
\end{quote}

Polite \textbf{\ja{こうかいします}}, \emph{kōkaishimasu}; a \ja{する} verb. Past (polite): \textbf{\ja{かれは} \ja{こうかいしました。}}, \emph{kare wa kōkai shimashita}.

\begin{cousinweb}
\textbf{\ja{こうかいする}} (\emph{kōkai suru}) is Sino-Japanese (kango) \emph{kōkai} with the native verb \emph{suru}, \textquotedblleft{}to do\textquotedblright{}: two Chinese characters read with their on'yomi, readings borrowed from Chinese. They are \emph{kō}, \textquotedblleft{}later\textquotedblright{} (Middle Chinese \emph{huwX}) and \emph{kai}, \textquotedblleft{}regret\textquotedblright{} (\emph{xwojH}) (KANJIDIC2; Middle Chinese in Baxter's transcription).
\end{cousinweb}

\begin{grammarlens}[title={What you've built}]
One more word for describing what you lived and felt.
\end{grammarlens}

\subsection*{Your turn}
\begin{itemize}
\raggedright
  \item \emph{Say it:} \emph{kōkai suru}
  \item \emph{Say it:} \emph{kōkai suru}, once more
  \item \emph{Say it:} \emph{ano toki no koto o kōkai shite imasu} — I regret what happened that time.
  \item \emph{Recall:} say the Japanese for to be troubled, then the Japanese for to brood, then say \emph{kōkai suru} again
\end{itemize}

\begin{tcolorbox}[breakable,colback=teal!4,colframe=teal!35!black,title={Before you move on}]
What does \ja{こうかいする} mean? (\textquotedblleft{}To regret\textquotedblright{}.) Say it once more.
\end{tcolorbox}
\section[\texorpdfstring{\ja{むなしい} (munashii)}{むなしい (munashii)}]{\ja{むなしい} (munashii) — empty; futile}

\label{lesson:JA-C432-munashii}



\begin{quote}
Before the new one: say the Japanese for to brood, then the Japanese for to regret.
\end{quote}

\subsection*{You'll want to know: \ja{むなしい}}
\textbf{\ja{むなしい}} — \emph{munashii} — \textquotedblleft{}empty; futile\textquotedblright{}.

\begin{quote}
\textbf{\ja{しあいに} \ja{まけて、むなしかったです。}} — \emph{shiai ni makete, munashikatta desu} — \textquotedblleft{}We lost the match, and it felt empty.\textquotedblright{}
\end{quote}

An i-adjective.

\begin{cousinweb}
\textbf{\ja{むなしい}} (\emph{munashii}) is a native Japanese word (wago), inherited rather than borrowed from Chinese; its kanji spelling was attached to it later (JMdict).
\end{cousinweb}

\begin{grammarlens}[title={What you've built}]
One more word for describing what you lived and felt.
\end{grammarlens}

\subsection*{Your turn}
\begin{itemize}
\raggedright
  \item \emph{Say it:} \emph{munashii}
  \item \emph{Say it:} \emph{munashii}, once more
  \item \emph{Say it:} \emph{shiai ni makete, munashikatta desu} — We lost the match, and it felt empty.
  \item \emph{Recall:} say the Japanese for to brood, then the Japanese for to regret, then say \emph{munashii} again
\end{itemize}

\begin{tcolorbox}[breakable,colback=teal!4,colframe=teal!35!black,title={Before you move on}]
What does \ja{むなしい} mean? (\textquotedblleft{}Empty; futile\textquotedblright{}.) Say it once more.
\end{tcolorbox}
\section[\texorpdfstring{\ja{あきらめる} (akirameru)}{あきらめる (akirameru)}]{\ja{あきらめる} (akirameru) — to give up; to resign oneself}

\label{lesson:JA-C432-akirameru}



\begin{quote}
Before the new one: say the Japanese for to regret, then the Japanese for futile.
\end{quote}

\subsection*{You'll want to know: \ja{あきらめる}}
\textbf{\ja{あきらめる}} — \emph{akirameru} — \textquotedblleft{}to give up; to resign oneself\textquotedblright{}.

\begin{quote}
\textbf{\ja{あめで} \ja{りょこうを} \ja{あきらめました。}} — \emph{ame de ryokō o akiramemashita} — \textquotedblleft{}We gave up on the trip because of the rain.\textquotedblright{}
\end{quote}

Polite \textbf{\ja{あきらめます}}, \emph{akiramemasu}; an ichidan verb. Past (polite): \textbf{\ja{せんしゅが} \ja{ゆめを} \ja{あきらめました。}}, \emph{senshu ga yume o akiramemashita}.

\begin{cousinweb}
\textbf{\ja{あきらめる}} (\emph{akirameru}) is a native Japanese word (wago), inherited rather than borrowed from Chinese; its kanji spelling was attached to it later (JMdict).
\end{cousinweb}

\begin{grammarlens}[title={What you've built}]
One more word for describing what you lived and felt.
\end{grammarlens}

\subsection*{Your turn}
\begin{itemize}
\raggedright
  \item \emph{Say it:} \emph{akirameru}
  \item \emph{Say it:} \emph{akirameru}, once more
  \item \emph{Say it:} \emph{ame de ryokō o akiramemashita} — We gave up on the trip because of the rain.
  \item \emph{Recall:} say the Japanese for to regret, then the Japanese for futile, then say \emph{akirameru} again
\end{itemize}

\begin{tcolorbox}[breakable,colback=teal!4,colframe=teal!35!black,title={Before you move on}]
What does \ja{あきらめる} mean? (\textquotedblleft{}To give up; to resign oneself\textquotedblright{}.) Say it once more.
\end{tcolorbox}
